\section{Introduction}
\label{sec:intro}

Transformers~\cite{vaswani2017attention} have emerged as the foundation of modern architectures for novel view synthesis~\cite{jinlvsm}, 3D reconstruction~\cite{jiang2025rayzer}, and camera-controllable video generation~\cite{bahmani2025ac3d,he2024cameractrl,cheong2024boosting}, where networks must reason about how visual observations are formed by camera geometries (\eg, pose, intrinsics, projection model, lens distortion) in order to ground pixel sequences into 3D space.
To maximize visual coverage from limited viewpoints, real-world applications in autonomous driving~\cite{yogamani2019woodscape,kumar2020fisheyedistancenet}, robotic perception~\cite{mei2006calibration,nayar1997catadioptric,zhang2016benefit}, and world models~\cite{team2025hunyuanworld,bar2025navigation,yang2025matrix} routinely use fisheye, catadioptric, and equirectangular projections for 360° panoramas.
However, despite progress in Transformer backbones~\cite{jinlvsm,jiang2025rayzer,wan2025wan,yang2025cogvideox}, most camera-conditioned methods still rely on simplified pinhole assumptions, overlooking the highly nonlinear projections and strong lens distortions present in these diverse camera geometries.

\begin{figure}[t]
    \centering
    \includegraphics[width=1.\linewidth, trim={0 0 0 0}, clip]{figure/pdf/rays.pdf}
    \vspace{-1.5em}
    \caption{
        \textbf{Comparison of camera encoding methods.}
        (a)~\emph{Direct Parameterization} encodes camera intrinsics and extrinsics as raw parameters, which lacks geometric interpretability and compatibility across camera types.
        (b)~\emph{Pl\"ucker Encoding} represents each ray as a pair of direction and moment vectors, providing a physically grounded but absolute, coordinate-dependent description.
        (c)~\emph{Projective Positional Encoding} encodes relative cameras in projective space, yet assumes pinhole projection and cannot model non-linear lens distortions.
        (d)~Our \emph{Relative Ray Encoding} reformulates geometric relationships in ray space, where each token corresponds to its own viewing ray, enabling better pose generalization and compatibility with \textbf{arbitrary} camera lenses.
    }\label{fig:rays}
    \vspace{-1em}
\end{figure}




In practice, these pinhole-based assumptions manifest in the camera representations fed into the network, as summarized in~\cref{fig:rays} and detailed in~\cref{sec:prelim}, which can be categorized into \emph{absolute} and \emph{relative} encodings.
Absolute encodings either feed raw camera parameters directly~\cite{bai2025recammaster,wang2024motionctrl} or adopt Pl\"ucker encoding~\cite{bahmani2025ac3d,he2024cameractrl} to represent rays under pinhole models.
However, as shown in our experiments (\cref{sec:cmp}), their dependence on specific world-frame configurations limits generalization across diverse camera geometries.
Recent relative encodings such as GTA~\cite{miyato2023gta} and PRoPE~\cite{li2025cameras} improve multi-view consistency in novel view synthesis by injecting geometry- or projection-aware relations into attention.
Yet, they remain restricted to the pinhole projection and are not readily compatible with pretrained video diffusion models.
\textit{Consequently, existing camera-conditioning methods still lack an encoding mechanism that explicitly models heterogeneous camera geometries, leading to limited generalization across varying intrinsics, distortion profiles, and viewpoint settings.}


To address these limitations, we propose \emph{Relative Ray Encoding} (\cref{fig:rays}d), which reformulates positional encoding from relative-camera encoding into relative-ray encoding.
By representing each token as a viewing ray expressed in local ray coordinates, this encoding enables the attention mechanism to operate directly in ray space rather than at the camera level, achieving geometry-consistent conditioning across heterogeneous camera lenses.

To evaluate this capability under diverse controllability demands, we adopt camera-controlled text-to-video generation as a testbed task.
However, existing camera-conditioned T2V methods~\cite{bahmani2025ac3d,he2024cameractrl} define camera poses only \emph{relative} to the first frame, leaving the global rotation frame ambiguous.
In particular, two degrees of freedom---\emph{pitch} and \emph{roll}---cannot be uniquely established, making it impossible to specify or reproduce the absolute orientation of the initial view.
To resolve this ambiguity, we introduce an \emph{Absolute Orientation Encoding} (third column of \cref{fig:teaser}), which anchors the camera to a gravity-aligned ``up'' direction.
This encoding explicitly provides absolute pitch and roll angles as a latitube-up map, enabling precise and repeatable control over the camera's absolute orientation.



Combining these two components, we arrive at \textbf{UCPE (Unified Camera Positional Encoding)}, a camera model-agnostic framework that injects complete camera geometry---including 6-DoF poses, intrinsics, and distortions---into Transformer attention.
UCPE provides a consistent formulation that unifies heterogeneous camera geometries, enabling fine-grained control over lens type, viewpoint, and orientation, as illustrated in~\cref{fig:teaser}.
Integrated into diffusion-based video generation through a lightweight spatial attention adapter, UCPE allows existing Diffusion Transformers to be fine-tuned with less than 1\% additional parameters.

To support systematic training and evaluation, we further construct a large-scale video dataset spanning diverse intrinsics, distortion profiles, and camera motions, offering a comprehensive benchmark for camera-controllable generation.
Experiments demonstrate that UCPE consistently improves camera controllability and visual fidelity across varied camera configurations, establishing a unified bridge between physical camera lenses and attention mechanisms for controllable video generation.

In summary, our contributions are as follows:
\begin{itemize}
    \item We propose \emph{UCPE}, a unified framework that encodes complete camera geometry (6-DoF poses, intrinsics, lens distortions) into Transformers. UCPE adopts a hybrid formulation combining \emph{Relative Ray Encoding} for ray-space geometric reasoning and projection nonlinearity and \emph{Absolute Orientation Encoding} for a gravity-aligned reference enabling controllable camera orientation.
    \item We introduce a lightweight \emph{attention injection method} that integrates UCPE into video Diffusion Transformers with less than 1\% of additional trainable parameters, enabling camera-aware video generation while preserving the visual fidelity of pretrained models.
    \item We construct a \emph{large-scale video dataset} covering diverse camera intrinsics, distortion profiles, and motion trajectories, providing a benchmark for evaluating camera-controllable generation under varied camera geometries.
    \item Extensive experiments demonstrate that UCPE not only enables accurate lens and orientation control, but also improves pose accuracy and generation quality for pinhole cameras and across diverse camera types.
\end{itemize}
